% LaTeX Canvas: Chapter Outline
\part{Computing Environment}
\chapter{Operating System}
\label{ch:os-management}

\section*{Purpose}
Your operating system (OS) is the platform that runs everything else: browsers, spreadsheets, programming tools, and file storage. Small OS maintenance mistakes—missed updates, broken permissions, full disks, weak account security, or missing backups—can derail assignments and corrupt data. This chapter gives novice students a practical, low-stress routine for keeping Windows and macOS stable, updated, and recoverable.

\section*{Learning objectives}
By the end of this chapter, a reader should be able to:
\begin{enumerate}
\item Find key OS settings on Windows and macOS using built-in search and navigation.
\item Check OS version/build information and interpret what matters for troubleshooting.
\item Understand types of updates and apply a safe patching workflow.
\item Set a baseline security posture (account hygiene, screen lock, encryption awareness).
\item Configure backups and understand restore/recovery options.
\item Apply a simple troubleshooting playbook (what to try first, what evidence to gather).
\item Establish a sustainable maintenance schedule for a semester.
\end{enumerate}

\section*{Running theme: stability comes from routines}
Most OS problems are predictable. The goal is not perfection; it is consistent, low-effort habits that prevent avoidable failures.

\section{A beginner mental model of the OS}
\subsection{What the OS is responsible for}
\begin{itemize}
\item User accounts, permissions, and security boundaries
\item Hardware and drivers (keyboard, Wi-Fi, storage)
\item Networking, storage, and application launching
\item Updates for the OS and built-in security protections
\item Recovery tools when something goes wrong
\end{itemize}

\subsection{What you control vs what you should avoid}
\begin{itemize}
\item You \textbf{should} manage settings, updates, backups, storage, and accounts.
\item You \textbf{should avoid} aggressive “cleaners,” registry hacks, random driver sites, or changing security settings without understanding the impact.
\end{itemize}

\section{Getting oriented: where settings live}
\subsection{Windows: Settings app, Windows Security, and the remaining Control Panel}
\begin{itemize}
\item Primary hub: \textbf{Settings} app (System, Windows Update, Apps, Accounts, Privacy \& security).
\item Security hub: \textbf{Windows Security} app (Defender, Firewall).
\item Legacy hub: \textbf{Control Panel} still hosts some tools (common example: File History).
\end{itemize}

\subsection{macOS: System Settings, Finder, and utilities}
\begin{itemize}
\item Primary hub: \textbf{System Settings} (General, Privacy \& Security, Network).
\item Utilities: \textbf{Disk Utility}, \textbf{System Information}, \textbf{Activity Monitor}.
\end{itemize}

\subsection{Settings search is your best friend}
\begin{itemize}
\item Use Settings/System Settings search to jump directly to the correct page.
\item Learn a small vocabulary: \texttt{update}, \texttt{backup}, \texttt{encryption}, \texttt{privacy}, \texttt{firewall}, \texttt{storage}, \texttt{recovery}.
\end{itemize}

\section{Know your system: version, storage, and constraints}
\subsection{Check OS version/build}
\begin{itemize}
\item Record OS version as part of your troubleshooting notes.
\item Understand the difference between:
\begin{itemize}
\item Major OS version (e.g., Windows 11; macOS \emph{name/version})
\item Feature releases/updates (Windows feature updates; macOS major upgrades)
\item Security/patch updates (smaller, frequent fixes)
\end{itemize}
\end{itemize}

\subsection{Check storage health (early and often)}
\begin{itemize}
\item Many failures are simply a full disk.
\item Learn to check free space and largest folders.
\item Keep a safety buffer of free space (rule of thumb: do not run at 99% full).
\end{itemize}

\subsection{Understand local vs synced storage}
\begin{itemize}
\item OneDrive (Windows) and iCloud Drive (macOS) can make files “on demand”.
\item For programming/data work, ensure key datasets are available locally before running.
\end{itemize}

\section{Updates and patching: best practices}
\subsection{Why updates matter}
\begin{itemize}
\item Security vulnerabilities are routinely discovered.
\item Updates also fix reliability issues and driver compatibility.
\end{itemize}

\subsection{Types of updates (student-relevant)}
\begin{description}
\item[Security updates] Fix vulnerabilities; highest priority.
\item[Quality/bug-fix updates] Improve stability; usually safe.
\item[Feature updates / major upgrades] Larger changes; can break workflows.
\item[Driver/firmware updates] Hardware compatibility; should come from OS/vendor.
\end{description}

\subsection{A safe patching workflow}
\begin{enumerate}
\item \textbf{Back up first} (or confirm backups are recent).
\item \textbf{Pick a good time window} (not 10 minutes before a deadline).
\item \textbf{Install updates} and allow required restarts.
\item \textbf{Verify} core tools still work (browser, spreadsheet, Python/conda, Jupyter).
\item \textbf{Record} what changed if something breaks (update history).
\end{enumerate}

\subsection{Windows update practices}
\begin{itemize}
\item Know where Windows Update settings live.
\item Understand restarts and “active hours”.
\item Use pause controls sparingly and intentionally.
\item Keep track of optional updates and driver updates.
\end{itemize}

\subsection{macOS update practices}
\begin{itemize}
\item Know where Software Update settings live.
\item Distinguish between point updates (safer) and major upgrades (plan ahead).
\item Consider enabling automatic updates, but still check periodically.
\end{itemize}

\section{Security baseline for students}
\subsection{Accounts and authentication}
\begin{itemize}
\item Use strong, unique passwords; prefer a password manager.
\item Enable multi-factor authentication (MFA) for your key accounts.
\item Avoid daily use of administrator accounts unless required.
\end{itemize}

\subsection{Screen lock and physical security}
\begin{itemize}
\item Require a password on wake.
\item Set a short auto-lock timer for laptops.
\end{itemize}

\subsection{Built-in protections}
\begin{itemize}
\item Windows: understand where to find Virus \& threat protection and Firewall settings.
\item macOS: understand privacy controls and app permission prompts.
\end{itemize}

\subsection{Encryption awareness (and recovery keys)}
\begin{itemize}
\item Encryption protects data if a device is lost or stolen, but it introduces a new dependency: \textbf{recovery keys}.
\item Windows: device encryption/BitLocker concepts; ensure you know where recovery keys are stored.
\item macOS: FileVault concepts; store recovery information safely.
\item Rule: never proceed with encryption changes unless you understand recovery.
\end{itemize}

\subsection{Software installation hygiene}
\begin{itemize}
\item Install software from official sources (vendor site, app store, trusted package managers).
\item Keep browsers updated.
\item Uninstall software you do not use.
\end{itemize}

\section{Backups and recovery: the student safety net}
\subsection{What backups are (and are not)}
\begin{itemize}
\item A backup is a second copy you can restore.
\item Sync is not always a backup (sync can propagate deletions).
\item Backups should be: automatic, regular, and restorable.
\end{itemize}

\subsection{Windows backups (conceptual map)}
\begin{itemize}
\item Windows Backup (settings and account-based restore workflows).
\item File History (versioned file copies; external drive/network location).
\end{itemize}

\subsection{macOS backups (conceptual map)}
\begin{itemize}
\item Time Machine (automatic backups to an external disk or approved location).
\item Practice restoring a file so you know it works.
\end{itemize}

\subsection{Recovery tools (when something breaks)}
\begin{itemize}
\item Windows: recovery options, reset, restore points (when enabled), recovery media.
\item macOS: macOS Recovery, Disk Utility First Aid, reinstall macOS.
\end{itemize}

\section{Performance and stability maintenance}
\subsection{Storage cleanup (do no harm)}
\begin{itemize}
\item Identify large files and old downloads.
\item Prefer deleting duplicates and stale installers.
\item Avoid deleting unknown system files.
\end{itemize}

\subsection{Startup and background apps}
\begin{itemize}
\item Too many startup items slow boot and drain battery.
\item Learn where to view and disable unnecessary login/startup apps.
\end{itemize}

\subsection{Battery and laptop health (practical notes)}
\begin{itemize}
\item Keep ventilation clear; avoid constant thermal throttling.
\item Replace failing batteries and chargers early if they create reliability issues.
\end{itemize}

\subsection{Drivers and peripherals}
\begin{itemize}
\item Prefer OS/vendor-provided updates.
\item Be cautious about third-party driver websites.
\item For printers and Wi-Fi, document the exact model and error messages.
\end{itemize}

\section{Troubleshooting playbook (student-friendly)}
\subsection{Start with the simplest reversible steps}
\begin{enumerate}
\item Restart the affected app.
\item Restart the computer.
\item Check free disk space.
\item Check network connectivity.
\item Identify what changed recently (updates, installs, new files).
\end{enumerate}

\subsection{Use safe modes and diagnostics}
\begin{itemize}
\item macOS: Safe Mode to isolate login items and caches.
\item Windows: safe boot/recovery tools (conceptually) when normal startup fails.
\end{itemize}

\subsection{Gather evidence for help}
\begin{itemize}
\item OS version/build, device model, and available disk space.
\item Exact error message and when it occurs.
\item Steps to reproduce.
\item Screenshot only when text cannot be copied.
\end{itemize}

\section{A semester maintenance schedule}
\subsection{Weekly (10 minutes)}
\begin{itemize}
\item Check for OS updates (or confirm automatic updates are active).
\item Confirm backups ran recently.
\item Quick storage check.
\end{itemize}

\subsection{Monthly (20–30 minutes)}
\begin{itemize}
\item Install pending feature/major updates only if you have time to recover.
\item Review startup apps.
\item Remove unused applications.
\end{itemize}

\subsection{Before major milestones (projects, finals)}
\begin{itemize}
\item Ensure backups are current and restorable.
\item Avoid major upgrades immediately before deadlines.
\item Confirm core tools (conda, notebook environment, editors) still run.
\end{itemize}

\section{Worked examples (outline)}
\subsection{Example 1: Set up a safe update routine}
\begin{itemize}
\item Identify Windows Update / macOS Software Update.
\item Enable/confirm automatic updates.
\item Schedule restarts to avoid deadline collisions.
\end{itemize}

\subsection{Example 2: Configure backups and test a restore}
\begin{itemize}
\item Windows: enable a backup approach and locate restore instructions.
\item macOS: set up Time Machine and restore a deleted file.
\end{itemize}

\subsection{Example 3: Recover from a broken update}
\begin{itemize}
\item Identify update history.
\item Use recovery options appropriately.
\item If needed: reinstall OS (last resort) after backing up.
\end{itemize}

\section{Exercises}
\begin{enumerate}
\item Record your OS version/build and device model in a troubleshooting note.
\item Find the update settings page and take a screenshot of where you would check for updates.
\item Configure a backup (Windows Backup/File History or Time Machine) and verify the last backup time.
\item Locate your encryption setting (Windows device encryption/BitLocker or macOS FileVault) and write down where the recovery key would be stored (do not paste the key into your notes).
\item Identify and disable one unnecessary startup/login item.
\item Simulate a recovery: restore one older version of a document from backup.
\end{enumerate}

\section{One-page checklist}
\begin{itemize}
\item I can find key settings (updates, backups, security, storage).
\item I know my OS version/build and can report it.
\item Updates are enabled and I install them on a schedule.
\item I avoid major upgrades right before deadlines.
\item Backups are configured, recent, and I have practiced a restore.
\item I keep adequate free disk space.
\item Screen lock is enabled and accounts use strong authentication.
\item I understand encryption recovery keys and where they are stored.
\item I can follow a basic troubleshooting playbook and gather evidence.
\end{itemize}

\appendix
\section{Quick map: where common settings live}
\subsection{Windows (typical)}
\begin{itemize}
\item \textbf{Windows Update:} Settings $\rightarrow$ Windows Update
\item \textbf{Security:} Windows Security app; Settings $\rightarrow$ Privacy \& security
\item \textbf{Recovery:} Settings $\rightarrow$ System $\rightarrow$ Recovery
\item \textbf{Backups:} Settings $\rightarrow$ Accounts/Backup areas; Control Panel for File History (on many systems)
\end{itemize}

\subsection{macOS (typical)}
\begin{itemize}
\item \textbf{Software Update:} System Settings $\rightarrow$ General $\rightarrow$ Software Update
\item \textbf{Privacy \& security:} System Settings $\rightarrow$ Privacy \& Security
\item \textbf{Encryption:} Privacy \& Security $\rightarrow$ FileVault
\item \textbf{Backups:} System Settings $\rightarrow$ General $\rightarrow$ Time Machine
\end{itemize}

% While \conda can have an intimidating level of complexity, your operating system is far more complex. This makes it all the more important to keep it up-to-date before starting to install \conda, Python, and everything else. 
An operating system manages all of the resources that you interact with while using a computer. For instance, it manages input from the keyboard, what is displayed on a screen, and how you open and close programs. When you interact with your computer, you are interacting with the operating system. For instance, Windows is an operating system. 

Keeping your operating system up-to-date ensures that you have the latest features, security, and bug fixes.

\section{I am confused about updating macOS.} If you have a Mac running macOS (previously called OS X), use the Software Update application to update macOS to the latest version.\footnote{More details and support can be found on the Apple support site \href{https://support.apple.com/en-us/HT201541}{here}.} Please note that this step could take 10--30 minutes, maybe longer.
\begin{marginfigure}
    \centering
    \includegraphics[width=\textwidth]{mac_software_update_app.png}
    \caption{The Windows Update application.}
    \label{fig:mac-update-app}
\end{marginfigure}
\begin{marginfigure}
    \centering
    \includegraphics[width=\textwidth]{mac_software_update_screen.png}
    \caption{Your screen while Windows is updating.}
    \label{fig:mac-update-screen}
\end{marginfigure}
\begin{enumerate}[noitemsep,topsep=0pt]
    \item Save all your open documents and close any applications you are not using, including documents and applications that may be minimized or on other spaces.
    \item In the upper-left corner of your screen, select the Apple menu 
\includegraphics[height=.5\baselineskip]{apple_logo.png} and choose ``System Preferences''.
    \item Select ``Software Update'' then ``Update Now'' or ``Upgrade Now''. See Figure~\ref{fig:mac-update-app}.
    \item Wait for the download to complete. This may take a few minutes or longer depending on the speed of your connection and the size of the update.
    \item Select the ``Restart Now'' button. Depending on your security settings, may need to enter your account password. You may be prompted to save and close other applications.
    \item Your computer should partially restart and you will see a black screen with a white Apple logo with a white status bar and time estimate below it. This bar may fill up and your computer may restart again one or more times, this is usually normal. See Figure~\ref{fig:mac-update-screen}.
    \item The update is complete when you are back at the login screen. Log in and check for updates one more time to be sure. If there are no more updates, proceed to the Step 1 section. If there are more updates, install them before proceeding.
\end{enumerate}

\section{I am confused about updating Windows.} If you have a PC running Windows, use the Windows Update application to update Windows to the latest version.\footnote{More details and support can be found on the Windows support site \href{https://support.microsoft.com/en-us/windows/update-windows-10-3c5ae7fc-9fb6-9af1-1984-b5e0412c556a}{here}.} Please note that this step could take 10--30 minutes, maybe longer.
\begin{marginfigure}
    \centering
    \includegraphics[width=\textwidth]{windows_update_app.png}
    \caption{The Windows Update application.}
    \label{fig:windows-update-app}
\end{marginfigure}
\begin{marginfigure}
    \centering
    \includegraphics[width=\textwidth]{windows_update_screen.png}
    \caption{Your screen while Windows is updating.}
    \label{fig:windows-update-screen}
\end{marginfigure}

\begin{enumerate}[noitemsep,topsep=0pt]
    \item Save all your open documents and close any applications you are not using, including documents and applications that may be minimized or on other desktops.
    \item Select the Start 
\includegraphics[height=.65\baselineskip]{windows_logo.png} button and go to ``Settings'', ``Update \& Security'', and select ``Windows Update''.
    \item An update may already be waiting. Click the ``Check for updates'' button to check if other updates are queued. See Figure~\ref{fig:windows-update-app}.
    \item Wait for the download to complete. This may take a few minutes or longer depending on the speed of your connection and the size of the update.
    \item Select the ``Restart Now'' button. Depending on your security settings, may need to enter your account password. You may be prompted to save and close other applications.
    \item Your computer should partially restart and you will see a blue screen with a progress indicator. See Figure~\ref{fig:windows-update-screen}.
    \item The update is complete when you are back at the login screen. Log in and check for updates one more time to be sure. If there are no more updates, proceed to the Step 1 section. If there are more updates, install them before proceeding.
\end{enumerate}

